\addtocontents{toc}{\protect\newpage}
\section{คู่มือการใช้งานสำหรับนิสิต}
\subsection{เข้าสู่ระบบและเมนู}
นิสิตเข้าสู่ระบบด้วยบัญชีที่มี \role{STUDENT} เมนูประกอบด้วย หน้าหลัก ใบลาของฉัน วิชาของฉัน ช่องทางการติดต่อ และรายงาน บนโทรศัพท์ให้ใช้เมนูสำหรับมือถือที่หน้าจอแสดง

\subsection{สร้างคำร้อง}
\begin{enumerate}
  \item เปิด \textbf{ใบลาของฉัน} และเลือกสร้างใบลาใหม่
  \item เลือก \textbf{ลาป่วย} หรือ \textbf{ลากิจ} แล้วกรอกเหตุผล
  \item เลือกวันเริ่มและสิ้นสุด โดยวันสิ้นสุดต้องไม่ก่อนวันเริ่ม
  \item ใช้เบอร์และที่อยู่จากระบบ หรือกรอกใหม่ให้ครบ
  \item แนบหลักฐาน ลาป่วยต้องแนบเสมอ
  \item เลือกรายวิชาที่ได้รับผลกระทบอย่างน้อยหนึ่งวิชา
  \item ตรวจ section อาจารย์ และระบุวันกับเวลาเรียนที่ลา
  \item ตรวจข้อมูล แล้วกดส่งหนึ่งครั้งและรอข้อความยืนยัน
\end{enumerate}
หลักฐานต้องเป็น JPG, JPEG, PNG, WEBP หรือ PDF ไม่เกิน 2 MB หากระบบไม่พบเส้นทางอนุมัติ ให้ติดต่อเจ้าหน้าที่พร้อมรหัสนิสิตและไม่ควรสร้างซ้ำ

\subsection{ติดตามและเปิดเอกสาร}
เปิด \textbf{ใบลาของฉัน} เลือกคำร้องเพื่อดูผู้พิจารณาปัจจุบันและผลแต่ละขั้น หากไม่ผ่านให้อ่านความคิดเห็น เมื่ออนุมัติแล้วสามารถเปิดเอกสารเพื่อพิมพ์ บันทึก PDF หรือใช้ QR/ลิงก์ตรวจสอบตามหน้าจอ

\subsection{รายวิชาและข้อมูลติดต่อ}
เมนู \textbf{วิชาของฉัน} ใช้ตรวจรายวิชาที่นำมาใช้ในแบบฟอร์ม เมนู \textbf{ช่องทางการติดต่อ} ใช้ตรวจหรือแก้เบอร์และที่อยู่ หากรายวิชาไม่ถูกต้องให้ตรวจภาคการศึกษาปัจจุบันและติดต่อเจ้าหน้าที่ก่อนส่งคำร้อง

เมื่อแจ้งปัญหา ควรเตรียมเลขคำร้อง รหัสนิสิต วันเวลา ขั้นตอน และข้อความผิดพลาด ห้ามส่ง token, รหัสผ่าน, cookie หรือหลักฐานทางการแพทย์ในช่องทางสาธารณะ
